\section{模型假设}

为使模型既符合题面，又避免把未给出的现实机制强行加入优化，作如下假设。
\begin{enumerate}[label=\arabic*., itemsep=2pt, topsep=3pt]
    \item \textbf{统一时间尺度。}以 $10$ min 为一个调度时段，一天共有 $T=144$ 个时段，$\Delta t=1/6$ h；附件中的 $0{:}00{+}1$ 视为当天 $24{:}00$。
    \item \textbf{不考虑反向售电。}题目只规定从外网购电，没有给出向外网售电的价格和交易规则，因此外网购电量均非负；无法被负载或储能吸收的光伏记为弃光。
    \item \textbf{储能参数在全年保持不变。}主模型取充电效率和放电效率均为 $0.9$，忽略自放电、容量衰减、循环寿命成本和启停损耗；效率解释的影响在灵敏度分析中另行讨论。
    \item \textbf{功率上限按时段折算。}最大充、放电功率 $5000$ kW 在 $10$ min 时段内对应最大充、放电量 $5000\Delta t=833.33$ kWh。
    \item \textbf{全年状态连续。}2025年1月1日 $0{:}00$ 储电量为 $6000$ kWh；随后日与日之间传递实际日末储电量。由于官方结果文件从2月1日开始，1月用于预测与储能状态预热。
    \item \textbf{问题二采用严格历史信息主口径。}题面未单独提供当天负载预报，故主模型只用决策时刻之前的实际负载构造因果预测；“日前负荷完全已知”仅作为理想信息基准，用于衡量信息价值。
    \item \textbf{冷启动使用附件1先验。}1月1日前没有历史观测时，以附件1的典型日负载、光伏曲线作为一次性先验；获得实际历史数据后即转为滚动历史预测。
    \item \textbf{允许储能进行实时偏差缓冲。}全天计划购电量锁定后，储能可以依据当前时段已经实现的负载、光伏和储电量修正充放电，但不得使用未来真实数据；储能仍无法覆盖的缺口才紧急购电。
    \item \textbf{问题二、三计划层采用日循环终端状态。}除问题一的题面硬约束外，问题二、三主模型令计划终端储电量等于当日实际初始储电量，以避免有限时域优化在日末透支电池；无终端约束情形作为敏感性对照。
    \item \textbf{问题四主结果直接使用附件4电价。}因果价格预测仅用于说明实际发布机制受限时可能产生的机会成本，不替代题面给定数据的正式重算结果。
\end{enumerate}

\section{符号说明}

为避免在后文反复解释，主要符号列于表~\ref{tab:symbols}。其中，功率量使用 kW，进入时段能量平衡前均乘以 $\Delta t$ 转为 kWh；$S_{d,t}$ 表示储能内部电量，而不是百分比形式的 SOC。

\begin{table}[htbp]
    \caption{主要符号说明}
    \label{tab:symbols}
    \centering
    \zihao{-5}
    \begin{tabularx}{0.94\textwidth}{>{\centering\arraybackslash}p{2.5cm} >{\centering\arraybackslash}p{1.8cm} X}
        \toprule
        \multicolumn{1}{c}{符号} & \multicolumn{1}{c}{单位} & \multicolumn{1}{c}{含义} \\
        \midrule
        $d,\ t$ & --- & 日期下标与日内 $10$ min 时段下标（$t=1,\dots,144$）\\
        $\Delta t$ & h & 单个调度时段长度，$\Delta t=1/6$ \\
        $L_{d,t}$ & kW & 小区负载功率 \\
        $P_{d,t}$ & kW & 光伏实际发电功率 \\
        $\hat P_{d,t}$ & kW & 决策时刻可用的光伏预测功率 \\
        $\pi_{d,t}$ & 元/kWh & 外网正常购电单价 \\
        $G^{p}_{d,t}$ & kWh & $0{:}00$ 制定的计划购电量 \\
        $G^{a}_{d,t}$ & kWh & 执行前最后一次有效调整后的购电量 \\
        $C_{d,t}$ & kWh & 母线侧进入储能的充电量 \\
        $H_{d,t}$ & kWh & 储能送往母线的放电量 \\
        $S_{d,t}$ & kWh & 第 $t$ 时段结束时储能内部储电量 \\
        $W_{d,t}$ & kWh & 未被利用的弃光电量 \\
        $E_{d,t}$ & kWh & 供电不足时的紧急购电量 \\
        $u_{d,t},\ v_{d,t}$ & kWh & 调整购电量相对计划的上调量 $(G^a-G^p)^+$ 与下调量 $(G^p-G^a)^+$ \\
        $z_{d,t}$ & --- & 充放电互斥变量；$z=1$ 允许充电，$z=0$ 允许放电 \\
        $\eta_c,\ \eta_d$ & --- & 充电效率与放电效率，主模型均取 $0.9$ \\
        $S_{\min},\ S_{\max}$ & kWh & 储电量上下限，分别为 $1200$ 与 $10800$ \\
        $\bar E$ & kWh & 单时段充放电量上限，$\bar E=5000\Delta t=833.33$ \\
        $\mathcal{I}_{d,\tau}$ & --- & 第 $d$ 天决策节点 $\tau$ 已经获得的信息集合 \\
        \bottomrule
    \end{tabularx}
\end{table}

表~\ref{tab:symbols} 中最容易混淆的是 $G^p$ 与 $G^a$：前者是 $0{:}00$ 对未来各时段作出的基准购电计划，后者是问题三中在日内新预报到达后、执行前最后一次有效调整后的购电量。二者的差额决定调整费用；实际运行仍可能因预测误差产生紧急购电 $E$。
